% ==============================================================
% Hypothesis Testing
% ==============================================================

\section{Hypothesis Testing}

\subsection{z-tests and t-tests}

% TODO

\subsection{p-values and critical regions}

% TODO

\subsection{Correlation and linear regression}

\subsubsection*{Ordinary Least Squares (OLS) Regression}

\begin{definition}
In \textbf{multiple linear regression}, we model a response variable $y$ as a linear function of $p$ predictors:
\[
\mathbf{y} = \mathbf{X}\boldsymbol{\beta} + \boldsymbol{\varepsilon}, \qquad \boldsymbol{\varepsilon} \sim (0, \sigma^2 \mathbf{I})
\]
where $\mathbf{y} \in \mathbb{R}^n$, $\mathbf{X} \in \mathbb{R}^{n \times p}$ is the \emph{design matrix}, and $\boldsymbol{\beta} \in \mathbb{R}^p$ is the coefficient vector. The \textbf{Ordinary Least Squares} (OLS) estimator minimises the sum of squared residuals:
\[
\hat{\boldsymbol{\beta}}_{\mathrm{OLS}} = \arg\min_{\boldsymbol{\beta}} \|\mathbf{y} - \mathbf{X}\boldsymbol{\beta}\|^2 = (\mathbf{X}^\top \mathbf{X})^{-1}\mathbf{X}^\top \mathbf{y}
\]
\end{definition}

\begin{remark}
The closed-form solution $(\mathbf{X}^\top \mathbf{X})^{-1}$ exists \emph{only if} $\mathbf{X}^\top\mathbf{X}$ is invertible, which requires $\mathbf{X}$ to have full column rank. This is precisely the condition violated by perfect multicollinearity.
\end{remark}

\subsubsection*{Derivation of the OLS Estimator}

The OLS estimator is found by minimising the \textbf{sum of squared residuals} (the loss function):
\[
\mathcal{L}(\boldsymbol{\beta}) = (\mathbf{y} - \mathbf{X}\boldsymbol{\beta})^\top(\mathbf{y} - \mathbf{X}\boldsymbol{\beta})
\]

\textbf{Step 1 — Expand the loss.} Using $(a-b)^\top(a-b) = a^\top a - 2b^\top a + b^\top b$:
\[
\mathcal{L}(\boldsymbol{\beta}) = \mathbf{y}^\top\mathbf{y} - 2\boldsymbol{\beta}^\top\mathbf{X}^\top\mathbf{y} + \boldsymbol{\beta}^\top\mathbf{X}^\top\mathbf{X}\boldsymbol{\beta}
\]

\textbf{Step 2 — Differentiate and set to zero.} Taking the gradient with respect to $\boldsymbol{\beta}$:
\[
\frac{\partial \mathcal{L}}{\partial \boldsymbol{\beta}} = -2\mathbf{X}^\top\mathbf{y} + 2\mathbf{X}^\top\mathbf{X}\boldsymbol{\beta} = \mathbf{0}
\]

\textbf{Step 3 — Normal equations.} Rearranging gives:
\[
\mathbf{X}^\top\mathbf{X}\,\boldsymbol{\beta} = \mathbf{X}^\top\mathbf{y}
\]

\textbf{Step 4 — Solve.} Provided $\mathbf{X}^\top\mathbf{X}$ is invertible, multiply both sides on the \emph{left} by $(\mathbf{X}^\top\mathbf{X})^{-1}$:
\[
\hat{\boldsymbol{\beta}}_{\mathrm{OLS}} = (\mathbf{X}^\top\mathbf{X})^{-1}\mathbf{X}^\top\mathbf{y}
\]

\begin{remark}
Note on order: writing $\mathbf{X}^\top\mathbf{y}\,(\mathbf{X}^\top\mathbf{X})^{-1}$ is \textbf{wrong} — matrix multiplication is not commutative. The inverse must appear on the left.
\end{remark}

\subsubsection*{The No Perfect Multicollinearity Assumption}

\begin{definition}
The \textbf{no perfect multicollinearity} assumption requires that no predictor variable is an exact linear combination of the others. Equivalently, the design matrix $\mathbf{X}$ must have full column rank: $\mathrm{rank}(\mathbf{X}) = p$.
\end{definition}

\begin{remark}
If perfect multicollinearity holds (e.g.\ $X_3 = 2X_1 + X_2$ exactly), then $\mathbf{X}^\top\mathbf{X}$ is singular and its inverse does not exist. OLS estimation breaks down completely --- there are infinitely many coefficient vectors that fit the data equally well, so no unique solution can be identified.
\end{remark}

\begin{example}
Suppose a model predicts income using both $X_1 = \text{age in months}$ and $X_2 = \text{age in years}$. Since $X_1 = 12 X_2$ exactly, the design matrix is rank-deficient. The model cannot separately estimate the effect of one month versus one year --- they always move together. \textbf{Fix:} drop one of the two variables.
\end{example}

\subsubsection*{Common Causes}

\begin{remark}
The three most frequent sources of perfect multicollinearity are:
\begin{itemize}
    \item \textbf{Dummy variable trap:} Including all $k$ categories of a categorical variable as dummy variables alongside an intercept. Since the dummies sum to 1 (the intercept column), perfect collinearity is introduced. \textbf{Fix:} drop one reference category.
    \item \textbf{Duplicate variables:} Including the same measurement twice in different units (e.g.\ area in $\text{m}^2$ and area in $\text{cm}^2$).
    \item \textbf{Constant predictors:} A predictor with zero variance (same value for all observations) is perfectly collinear with the intercept.
\end{itemize}
\end{remark}

\subsubsection*{Perfect vs.\ Imperfect Multicollinearity}

\begin{definition}
\textbf{Imperfect (near) multicollinearity} occurs when predictors are highly correlated but not in an exact linear relationship. OLS still produces a unique solution, but coefficient estimates become highly sensitive to small changes in the data and standard errors inflate, making inference unreliable.
\end{definition}

\begin{remark}
The \textbf{Variance Inflation Factor} (VIF) measures the severity of imperfect multicollinearity for predictor $j$:
\[
\mathrm{VIF}_j = \frac{1}{1 - R_j^2}
\]
where $R_j^2$ is the $R^2$ from regressing $X_j$ on all other predictors. A rule of thumb: $\mathrm{VIF}_j > 10$ signals problematic multicollinearity.
\end{remark}

\subsubsection*{Heteroscedasticity}

\begin{definition}
The OLS error term $\varepsilon_i$ is \textbf{homoscedastic} if its conditional variance is constant across all observations:
\[
\mathrm{Var}(\varepsilon_i \mid X_i) = \sigma^2 \quad \forall\, i
\]
It is \textbf{heteroscedastic} if this variance instead depends on $X_i$:
\[
\mathrm{Var}(\varepsilon_i \mid X_i) = \sigma_i^2 \quad \text{(varies with } i\text{)}
\]
\end{definition}

\textbf{Variance structure.} Common parametric forms for $\sigma_i^2$ include:
\[
\sigma_i^2 = \sigma^2 \cdot X_i, \qquad \sigma_i^2 = \sigma^2 \cdot X_i^2, \qquad \sigma_i^2 = \exp(\alpha + \beta X_i)
\]
The subscript $i$ is the essential signal: each observation carries its own variance rather than a single pooled value.

\textbf{Visual signature.} A residual-vs-fitted plot showing a \emph{funnel} (spread growing with fitted values) is the classic indicator of heteroscedasticity. A reverse funnel is equally diagnostic.

\begin{remark}[Consequences for OLS]
Under heteroscedasticity:
\begin{itemize}
    \item \textbf{Coefficients remain unbiased} — $\hat{\boldsymbol{\beta}}_{\mathrm{OLS}}$ still converges to the true $\boldsymbol{\beta}$ in expectation.
    \item \textbf{Gauss--Markov fails} — OLS is no longer the Best Linear Unbiased Estimator (BLUE); more efficient estimators (e.g.\ WLS) exist.
    \item \textbf{Standard errors are biased} — standard errors computed under homoscedasticity are typically \emph{underestimated}, inflating $t$-statistics.
    \item \textbf{Inference breaks down} — because standard errors feed into $t$-tests, $F$-tests, and confidence intervals, all classical hypothesis tests become unreliable and may produce spurious rejections of $H_0$.
\end{itemize}
\end{remark}

\textbf{Detection tests.}
\begin{itemize}
    \item \textbf{Visual:} Plot residuals vs.\ fitted values; look for funnel or fan shapes.
    \item \textbf{Breusch--Pagan test:} Regress squared residuals $\hat{\varepsilon}_i^2$ on the regressors. The LM test statistic follows $\chi^2(k)$ under $H_0$ (homoscedasticity).
    \item \textbf{White test:} Like Breusch--Pagan but also includes squares and cross-products of regressors, making it more robust to functional-form misspecification.
    \item \textbf{Goldfeld--Quandt test:} Split the sample into two groups; compare their residual variances via an $F$-test.
\end{itemize}

\begin{remark}[Remedies]
\begin{itemize}
    \item \textbf{Heteroscedasticity-consistent (HC) standard errors} (robust SEs): correct standard error estimates without assuming constant variance. Common variants are HC0 (White 1980), HC1 (degrees-of-freedom corrected), and HC3 (recommended for small samples).
    \item \textbf{Weighted Least Squares (WLS):} down-weights high-variance observations. Efficient when the variance structure $\sigma_i^2$ is known or estimable.
    \item \textbf{Generalised Least Squares (GLS):} full generalisation of WLS for an arbitrary variance--covariance structure.
    \item \textbf{Variance-stabilising transformations:} apply $\ln(Y)$ or a power transformation to the dependent variable to stabilise the spread of residuals.
\end{itemize}
\end{remark}

\subsubsection*{Joint Hypothesis Testing: the $F$-Statistic in Regression}

A \textbf{$t$-test} checks one coefficient at a time. An \textbf{$F$-test} checks several coefficients simultaneously. To test $H_0\colon \beta_1 = 0$ \emph{and} $\beta_2 = 0$ together, combine both $t$-statistics into a single score.

\begin{theorem}[Joint $F$-statistic, two restrictions]
Let $t_1 = \hat\beta_1 / \widehat{\mathrm{SE}}_1$ and $t_2 = \hat\beta_2 / \widehat{\mathrm{SE}}_2$ be the individual $t$-statistics, and let $\hat\rho \equiv \hat\rho_{t_1,t_2}$ be the estimated correlation between them. Then the joint $F$-statistic is
\[
F = \frac{1}{2} \cdot \frac{t_1^2 + t_2^2 - 2\hat\rho\, t_1 t_2}{1 - \hat\rho^2}.
\]
Under $H_0$, in large samples, $F \sim F(2,\infty)$.
\end{theorem}

\textbf{Anatomy of the formula.}
\begin{itemize}
    \item $t_1^2 + t_2^2$: raw evidence. Squaring ensures large negative and large positive $t$-values both count.
    \item $-2\hat\rho\, t_1 t_2$: removes double-counting. If $t_1$ and $t_2$ move together ($\hat\rho > 0$), they share information; this term subtracts the overlap.
    \item $1 - \hat\rho^2$: rescaling. Even after removing the overlap in the numerator, correlated variables do not spread like independent ones — this denominator adjusts for the ``compressed'' joint geometry.
    \item $\frac{1}{2}$: divides by the number of restrictions (2), making the result comparable to $F(2,\cdot)$.
\end{itemize}

\begin{remark}[Special case: uncorrelated $t$-statistics]
When $\hat\rho = 0$ (regressors uncorrelated), the formula collapses to
\[
F = \tfrac{1}{2}(t_1^2 + t_2^2).
\]
This is simply the average of the squared $t$-statistics. In large samples, $t_1, t_2 \approx N(0,1)$ independently under $H_0$, so $t_1^2 + t_2^2 \sim \chi^2(2)$ and $F \sim F(2,\infty)$.
\end{remark}

\begin{remark}[General: $q$ restrictions]
Testing $q$ restrictions simultaneously gives $F \sim F(q,\infty)$ in large samples. The $q$ in the numerator degrees of freedom is precisely the number of restrictions being tested.
\end{remark}

\subsubsection*{Mahalanobis Distance and the $F$-Formula}

The joint $F$-statistic is a scaled \textbf{Mahalanobis distance} — the correct notion of distance when two variables are correlated.

Under $H_0$, $(t_1, t_2)^\top$ has covariance matrix
\[
\Sigma = \begin{pmatrix} 1 & \rho \\ \rho & 1 \end{pmatrix},
\]
where each $t$-stat has variance 1 (both are approximately $N(0,1)$) and they share correlation $\rho$ because $\hat\beta_1, \hat\beta_2$ come from the same regression data.

Inverting a $2\times 2$ matrix $\begin{pmatrix}a&b\\c&d\end{pmatrix}$ via $\frac{1}{ad-bc}\begin{pmatrix}d&-b\\-c&a\end{pmatrix}$, with $\det(\Sigma) = 1 - \rho^2$:
\[
\Sigma^{-1} = \frac{1}{1-\rho^2}\begin{pmatrix} 1 & -\rho \\ -\rho & 1 \end{pmatrix}.
\]

The Mahalanobis squared distance from the origin is $\mathbf{t}^\top \Sigma^{-1} \mathbf{t}$:
\[
\mathbf{t}^\top \Sigma^{-1} \mathbf{t}
= \frac{1}{1-\rho^2}(t_1, t_2)\begin{pmatrix}1 & -\rho \\ -\rho & 1\end{pmatrix}\begin{pmatrix}t_1\\t_2\end{pmatrix}
= \frac{t_1^2 + t_2^2 - 2\rho t_1 t_2}{1-\rho^2}.
\]
The $-2\rho t_1 t_2$ arises because the off-diagonal $-\rho$ entries appear in both rows of the matrix (symmetry), contributing $-\rho t_1 t_2$ twice. The $1-\rho^2$ is the determinant of $\Sigma$, falling out of the matrix inversion.

Dividing by 2 (the number of restrictions) gives the $F$-statistic $F = \frac{1}{2}\mathbf{t}^\top \hat\Sigma^{-1}\mathbf{t}$.

\begin{remark}[Geometric intuition: circles vs.\ ellipses]
When $\rho = 0$, the equal-distance contours of $(t_1,t_2)$ from the origin are circles. When $\rho \neq 0$, the cloud of possible values stretches along the diagonal — equal-distance contours become ellipses. Points along $t_1 \approx t_2$ are not genuinely independent evidence; they lie in a correlated direction. The Mahalanobis formula measures distance in this elliptical geometry, not the naive circular one. As $|\rho| \to 1$, the ellipse degenerates and the test becomes unstable ($1-\rho^2 \to 0$).
\end{remark}

\begin{remark}[Why $t$-statistics can be correlated in regression]
The $t$-stats $t_j = \hat\beta_j / \widehat{\mathrm{SE}}_j$ share the same sample data. If regressors $X_1$ and $X_2$ are correlated with each other, the coefficient estimates $\hat\beta_1$ and $\hat\beta_2$ are correlated too — they trade off against each other to fit the same outcomes. This is why $\hat\rho \neq 0$ is the norm, not the exception.
\end{remark}
